% Librerias
\usepackage{babel}
\babelprovide[main,import]{spanish}
\usepackage[utf8]{inputenc}


\usepackage{amsmath, amssymb}
\usepackage{siunitx}
\sisetup{
	output-decimal-marker = {,},  % Coma decimal
	inter-unit-product = \cdot,   % Usa \cdot entre unidades
	exponent-product = \cdot,	  % Usa \cdot entre mantisa y exponente
	per-mode = symbol             % Formato de unidades (/)
}

\usepackage{tabularx}
\usepackage{booktabs}
\newcolumntype{Y}{>{\raggedright\arraybackslash}X}
\newcolumntype{L}[1]{>{\raggedright\arraybackslash}p{#1}}
\usepackage{float}

\usepackage{graphicx}
\usepackage{caption}
\captionsetup[figure]{justification=centering, singlelinecheck=false}
\captionsetup[table]{justification=centering, singlelinecheck=false}

\usepackage[top=2.5cm, bottom=2.5cm, left=1.5cm, right=1.5cm]{geometry}
\raggedbottom
\setlength{\emergencystretch}{2em}

\usepackage{fancyhdr}
\pagestyle{fancy}
\setlength{\headheight}{14.5pt}

\usepackage[colorlinks=true, allcolors=black]{hyperref}
\fancyhead{}
\fancyhead[L]{UNC - FCEFyN} % Alineado a la izquierda
\fancyhead[R]{Redes de Computadoras} % Alineado a la derecha
\fancyfoot{}
\fancyfoot[C]{\thepage}

% Reduce el espacio vertical excesivo que la clase book agrega antes de cada capitulo.
\makeatletter
\def\@makechapterhead#1{%
  \vspace*{-50\p@}%
  {\parindent \z@ \raggedright \normalfont
    \ifnum \c@secnumdepth >\m@ne
      \Large\bfseries \@chapapp\space \thechapter
      \par\nobreak
      \vskip 6\p@
    \fi
    \huge \bfseries #1\par\nobreak
    \vskip 10\p@
  }}
\def\@makeschapterhead#1{%
  \vspace*{-50\p@}%
  {\parindent \z@ \raggedright
    \normalfont
    \huge \bfseries #1\par\nobreak
    \vskip 10\p@
  }}
\makeatother
